% Part: applied-modal-logics
% Chapter: epistemic-logic
% Section: properties-accessibility

\documentclass[../../../include/open-logic-section]{subfiles}

\begin{document}

\olfileid{aml}{el}{acc}

\olsection{روابط دسترسی‌پذیری و اصول معرفتی}

با توجه به آنچه از تناظر قاب در منطق‌های موجهات عادی می‌دانیم، شاید
بخواهیم ببینیم !!{formula}های مشخصه با تفسیرهای معرفتی چه صورتی پیدا
می‌کنند. پیش‌تر گفتیم که منطق‌های معرفتی را معمولاً در~S5 تفسیر
می‌کنند. پس ببینیم اصول گوناگون معرفتی چگونه بازنمایی می‌شوند و چگونه
با شرایط مختلف قاب تناظر دارند.

به یاد آورید که در منطق موجهات عادی، !!{formula}های موجهاتی گوناگون
خواص متفاوت روابط دسترسی‌پذیری را مشخص می‌کردند. جدول زیر چند مورد را
برمی‌گزیند که با اصول معرفتی معینی متناظرند.

\begin{table}[t]
    \begin{tabular}{| p{.48\textwidth} || p{.48\textwidth} |}
      \hline
      {\emph{اگر~$R$ \dots{} باشد}} & {\emph{آنگاه \dots{} در~$\mModel{M}$ صادق است:}} \\
      \hline \hline
        
      & $\Knows (p \lif q) \lif (\Knows p \lif \Knows q)$
      \hfill \newline{(بستار)} \\
      \hline
      \emph{بازتابی}: $\forall w Rww$
      & $\Knows p \lif p$ \hfill \newline{(صدق‌مندی)} \\
      \hline
      \emph{تعدی}: \newline
      $\forall u \forall v \forall w ((Ruv \land Rvw) \lif Ruw)$ &
      $\Knows p \lif \Knows \Knows p$ \hfill
           \newline{(درون‌نگری مثبت)} \\
      \hline
      \emph{اقلیدسی}: \newline
      $\forall w \forall u \forall v ((Rwu \land Rwv) \lif Ruv)$ &
      $\lnot \Knows p \lif \Knows \neg \Knows p$ \hfill
        \newline{(درون‌نگری منفی)} \\
      \hline
    \end{tabular}
    \caption{چهار اصل معرفتی.}
    \ollabel{tab:four}
  \end{table}

صدق‌مندی، که با اصل موضوع~$T$ متناظر است، اغلب کم‌مناقشه‌ترینِ این اصول
دانسته می‌شود، زیرا بیان می‌کند که اگر !!a{formula} دانسته باشد، باید
صادق باشد. بستار و نیز درون‌نگری مثبت و منفی بسیار مناقشه‌برانگیزترند.

بستار، که با اصل موضوع~$K$ متناظر است، این اندیشه را بازمی‌نماید که
دانش عامل تحت استلزام بسته است. این امر ممکن است در برخی موارد
پذیرفتنی بنماید. برای نمونه، شاید بدانم که اگر در ویکتوریا باشم، در
جزیرهٔ ونکوورم. اگر سناریوهای شکاکانهٔ نامتعارف را کنار بگذاریم، می‌دانم
که در ویکتوریا هستم، و ازاین‌رو باید بدانم که در جزیرهٔ ونکوورم. پس در
این مورد، بستار منطقیِ دانشم نسبتاً شهودی می‌نماید. از سوی دیگر، ما
همیشه همهٔ پیامدهای دانش خود را تا انتها بررسی نمی‌کنیم و بنابراین این
اصل در موارد دیگر ممکن است به نتایجی کم‌تر شهودی بینجامد.

درون‌نگری مثبت، که گاه اصل~KK نامیده می‌شود، گاهی چنین صورت‌بندی
می‌گردد: اگر چیزی را بدانم، آنگاه می‌دانم که آن را می‌دانم. این اصل
همتای معرفتیِ اصل موضوع~4 است. به همین ترتیب، درون‌نگری منفی چنین
صورت‌بندی می‌شود: اگر چیزی را \emph{ندانم}، آنگاه می‌دانم که آن را
نمی‌دانم؛ و این همتای اصل موضوع~5 است. برای هر دوی این اصول به نظر
می‌رسد مثال‌های نقضی نسبتاً عادی وجود دارد که در آن‌ها دربارهٔ اینکه
آیا چیزی را می‌دانم یا نه نامطمئنم، هرچند در واقع آن را می‌دانم.

\end{document}
